% ej_coord_R2_01.tex
%
% Copyright (C) 2022--2026 José A. Navarro Ramón <janr.devel@gmail.com>
% Licencia del código GPLv2
% Licencia Creative Commons Recognition Non-Commercial Share-alike.
% (CC-BY-NC-SA)

\section{Ejercicio 1}
Supongamos un cierto vector en coordenadas rectangulares,
$\vvv{v} = 3\,\uvec{\i} + 4\,\uvec{\j}$.\\
a) Calcula el cuadrado de su módulo.\\
b) Nos proponemos realizar un cambio de base desde el sistema rectangular a otro oblicuo,
con un ángulo entre ejes de \qty{60}{\degree}. Calcula la expresión de los vectores
unitarios de la base nueva. ¿Cuál es su módulo?
Escribe la matriz $\mmm{M}$ de cambio de base y su inversa.\\
c) Obtén el tensor métrico en el sistema oblicuo anterior.\\
d) Halla las componentes contravariantes del vector.\\
e) ¿Cuáles son las componentes covariantes?\\
f) Halla el cuadrado del módulo del vector en el sistema oblicuo.\\

\noindent
\textbf{Solución}

\noindent
a) El cuadrado del módulo del vector es
\[
  v^2 = v_x^2 + v_y^2 = 3^2 + 4^2 = 25
\]

\noindent
b) Utilizando la figura \ref{fig:R2-coord_oblicuas_e1e2}, la nueva base estaría formada
por dos vectores unitarios, que en función del sistema rectangular son
\begin{align*}
  \vvv{e}_1 &= \uvec{\i}
  \longrightarrow |\vvv{e}_1| = 1\\ 
  \vvv{e}_2 &= \cos\pi/3\,\uvec{\i} + \sin\pi/3\,\uvec{\j}
               = \dfrac{1}{2}\,\uvec{\i} + \dfrac{\sqrt{3}}{2}\,\uvec{\j}
               \longrightarrow |\vvv{e}_2| = \cos^2\pi/3 + \sin^2\pi/3 = 1
\end{align*}

La nueva base sería
\[
  \tilde{B} = \set{\vvv{e}_1, \vvv{e}_2}
  = \Set{\uvec{\i}, \dfrac{1}{2}\,\uvec{\i} + \dfrac{\sqrt{3}}{2}\,\uvec{\j}}
  = \Set{(1,0),\, (1/2, \sqrt{3}/2)}
\]

De acuerdo con esto, la matriz de cambio de base estaría formada por los vectores de la
nueva base en función de la antigua, dispuestos en columnas
\[
  \mmm{M}
  = \begin{pmatrix}
    1 & 1/2\\
    0 & \sqrt{3}/2
  \end{pmatrix}
\]

{\footnotesize
Comprobación
\[
  \tilde{B} = B \mmm{M}
\]

Los vectores de la base los disponemos en filas
\[
  \begin{pmatrix}
    \vvv{e}_1 & \vvv{e}_2
  \end{pmatrix}
  =
  \begin{pmatrix}
    \uvec{\i} & \uvec{\j}
  \end{pmatrix}
  \,
  \begin{pmatrix}
    1 & 1/2\\
    0 & \sqrt{3}/2
  \end{pmatrix}
  =
  \begin{pmatrix}
    \uvec{\i} & \dfrac{1}{2}\,\uvec{\i} + \dfrac{\sqrt{3}}{2}\,\uvec{\j}
  \end{pmatrix}
\]
}
La inversa de la matriz de cambio de base se calcula multiplicando la inversa del
determinante de la matriz por la adjunta de su transpuesta. El resultado es
\[
  \mmm{M}^{-1}
  = \dfrac{1}{\det(\mmm{M})}\,\text{adj}\left(\mmm{M}^\transp\right)
  = \dfrac{1}{\sqrt{3}/2}
  \begin{pmatrix}
    \sqrt{3}/2 & -1/2\\[1ex]
    0 & 1
  \end{pmatrix}
  =
  \begin{pmatrix}
    1 & -\sqrt{3}/3\\[1ex]
    0 & 2\sqrt{3}/3
  \end{pmatrix}
\]

\noindent
c) El tensor métrico en el sistema oblicuo, con $\alpha = \pi/3\,\unit{\radian}$ es
\[
  \mmm{g}_{ij}
  =\begin{pmatrix}
    g_{11} & g_{12}\\
    g_{21} & g_{22}
  \end{pmatrix}
  = \begin{pmatrix}
    \vvv{e}_1\cdot\vvv{e}_1 & \vvv{e}_1\cdot\vvv{e}_2\\
    \vvv{e}_2\cdot\vvv{e}_1 & \vvv{e}_2\cdot\vvv{e}_2\\
  \end{pmatrix}
  = \begin{pmatrix}
    1 & \cos\pi/3\\
    \cos\pi/3 & 1
  \end{pmatrix}
  =
  \begin{pmatrix}
    1 & 1/2\\
    1/2 & 1
  \end{pmatrix}
\]

Su inversa es
\[
  \mmm{g}^{ij}
  = \dfrac{1}{\det(\mmm{g}_{ij})}\, \text{adj}(\mmm{g}_{ij}^\transp)
  = \dfrac{1}{3/4}
  \begin{pmatrix}
    1 & -1/2\\
    -1/2 & 1
  \end{pmatrix}
  = \begin{pmatrix}
    4/3 & -2/3\\
    -2/3 & 4/3
  \end{pmatrix}
\]

\noindent
d) Para hallar las componentes contravariantes, debemos aplicar la inversa de la matriz
de cambio de base
\[
  \left[\tilde{x}^i\right]_{\tilde{B}} = \mmm{M}^{-1}\, \left[x^i\right]_{B}
\]
donde las componentes contravarianes en la base rectangular $[x^i]_{B}$ son iguales
que las covariantes y valen $(3, 4)$.

\[
  \begin{pmatrix}
    x^1\\
    x^2
  \end{pmatrix}
  = \begin{pmatrix}
    1 & -\dfrac{\sqrt{3}}{3}\\[1.5ex]
    0 & \dfrac{2\sqrt{3}}{3}
  \end{pmatrix}
  \,
  \begin{pmatrix}
    3\\
    4
  \end{pmatrix}
  = \begin{pmatrix}
    3 - \dfrac{4\sqrt{3}}{3}\\[1.5ex]
    \dfrac{8\sqrt{3}}{3}
  \end{pmatrix}
\]

\noindent
e) Las componentes covariantes se obtienen bajando índices
\[
  x_i = g_{ij}\,x^j
\]
Desarrollando la expresión para $x_1$ y $x_2$ y teniendo en cuenta la métrica,
$g_{11} = 1$, $g_{12}=g_{21}=1/2$ y $g_{22}=1$
\begin{align*}
  x_1
  &= g_{11}x^1 + g_{12}x^2 = x^1 + \dfrac{1}{2}x^2
    = 3 - \cancelout{\dfrac{4\sqrt{3}}{3}} + \cancelout{\dfrac{1}{2}\dfrac{8\sqrt{3}}{3}}
    = 3\\
  x_2
  &= g_{21}x^1 + g_{22} x^2 = \dfrac{1}{2}x^1 + x^2
    = \dfrac{1}{2}\left(3-\dfrac{4\sqrt{3}}{3}\right) + \dfrac{8\sqrt{3}}{3}
    = \dfrac{3}{2} + 2\sqrt{3}
\end{align*}

\noindent
f) El cuadrado del módulo del vector calculado con las componentes oblicuas debe ser
$25$, pues tiene que dar en mismo resultado que en rectangulares (apartado a).
\begin{align*}
  \vvv{v}\cdot\vvv{v}
  &=
    x^ix_i
    = x^1 x_1 + x^2 x_2
    = \left(3-\dfrac{4\sqrt{3}}{3}\right) 3
    + \dfrac{8\sqrt{3}}{3} \left(\dfrac{3}{2} + 2\sqrt{3}\right)\\
  &=
    9 - 4\sqrt{3}
    + \dfrac{\cancelout{3}\cdot 8\sqrt{3}}{2\cdot\cancelout{3}}
    + \dfrac{2\cdot 8\cdot \cancelout{3}}{\cancelout{3}}
    = 9 - \cancelout{4\sqrt{3}} + \cancelout{4\sqrt{3}} + 16
    = 25
\end{align*}



%%% Local Variables:
%%% mode: latex
%%% TeX-engine: luatex
%%% TeX-master: "../retazosmatematicas.tex"
%%% End:
